% Options for packages loaded elsewhere
\PassOptionsToPackage{unicode}{hyperref}
\PassOptionsToPackage{hyphens}{url}
%
\documentclass[
  10pt,
]{article}
\usepackage{lmodern}
\usepackage{amssymb,amsmath}
\usepackage{ifxetex,ifluatex}
\ifnum 0\ifxetex 1\fi\ifluatex 1\fi=0 % if pdftex
  \usepackage[T1]{fontenc}
  \usepackage[utf8]{inputenc}
  \usepackage{textcomp} % provide euro and other symbols
\else % if luatex or xetex
  \usepackage{unicode-math}
  \defaultfontfeatures{Scale=MatchLowercase}
  \defaultfontfeatures[\rmfamily]{Ligatures=TeX,Scale=1}
\fi
% Use upquote if available, for straight quotes in verbatim environments
\IfFileExists{upquote.sty}{\usepackage{upquote}}{}
\IfFileExists{microtype.sty}{% use microtype if available
  \usepackage[]{microtype}
  \UseMicrotypeSet[protrusion]{basicmath} % disable protrusion for tt fonts
}{}
\makeatletter
\@ifundefined{KOMAClassName}{% if non-KOMA class
  \IfFileExists{parskip.sty}{%
    \usepackage{parskip}
  }{% else
    \setlength{\parindent}{0pt}
    \setlength{\parskip}{6pt plus 2pt minus 1pt}}
}{% if KOMA class
  \KOMAoptions{parskip=half}}
\makeatother
\usepackage{xcolor}
\IfFileExists{xurl.sty}{\usepackage{xurl}}{} % add URL line breaks if available
\IfFileExists{bookmark.sty}{\usepackage{bookmark}}{\usepackage{hyperref}}
\hypersetup{
  hidelinks,
  pdfcreator={LaTeX via pandoc}}
\urlstyle{same} % disable monospaced font for URLs
\usepackage[a4paper,margin=18mm]{geometry}
\setlength{\emergencystretch}{3em} % prevent overfull lines
\providecommand{\tightlist}{%
  \setlength{\itemsep}{0pt}\setlength{\parskip}{0pt}}
\setcounter{secnumdepth}{-\maxdimen} % remove section numbering

\usepackage{esint}
\usepackage{xeCJK}
\setmainfont{DejaVu Serif}
\setmonofont{DejaVu Sans Mono}
\setCJKmainfont{Noto Serif CJK SC}
\setCJKmonofont{Noto Sans Mono CJK SC}
\setmathfont{Asana Math}
\setmathfont{Latin Modern Math}[range=\lrcorner]
\setlength{\emergencystretch}{3em}
\allowdisplaybreaks
\hypersetup{hypertexnames=false}

\author{}
\date{}

\begin{document}

\hypertarget{signed-quartic-completion-of-the-actual-torsion-flux}{%
\section{Signed quartic completion of the actual torsion
flux}\label{signed-quartic-completion-of-the-actual-torsion-flux}}

28 September 2026. Continuation of I41, Section 7, I51, Section 2, and
the \href{20260928_actual_torsion_prolongation.md}{actual torsion
prolongation}. The calculation below reorganizes the complete global
expression, rather than dropping the symmetric Hessian term in a
pointwise commutator. It yields an exact first-derivative identity and a
quantitative necessary condition for an endpoint sequence. The final
subthreshold estimate is not derived from the original hypotheses; O10
remains open.

\hypertarget{the-actual-field-and-the-characteristic-identity}{%
\subsection{1. The actual field and the characteristic
identity}\label{the-actual-field-and-the-characteristic-identity}}

Use the scalar-minus-four, constant-weight actual solution
\(D_A\Phi=0,\ F_A^+=2i\sigma(\Phi)\) of I51. Let tau be its cubic
opposite-spinor torsion, a section of \(E=S^-\otimes L\), and put
\(m=|\tau|^2,\ m_*=128/81\). Fix \(0<a<b<m_*\) and a smooth
nondecreasing function vartheta\_0, zero on {[}0,a{]} and one on
{[}b,infinity). Write

\[
 \vartheta=\vartheta_0(m),\quad \chi=1-\vartheta^2,\quad
                         T=\{a<m<b\}.
\]

On m \textgreater{} 0, use a local unitary frame with first vector u =
tau/sqrt(m). Write the actual connection as

\[
 \Gamma=\begin{pmatrix}\alpha&-\bar b_1\\b_1&\alpha_2\end{pmatrix},
 \qquad \alpha=i\theta,\quad b_1=\gamma+i\delta,
\]

where theta,gamma,delta are real one-forms. The phase one-form theta is
not the scalar cutoff vartheta. All expressions with cutoff derivatives
are supported in T, bounded away from tau's zero set.

In the inherited curvature and orientation conventions, the rank-two
Chern transgression is

\[
 4\pi^2\mathcal T_\tau
 =-\alpha\wedge F_{22}
 +\tfrac12(b_1\wedge F_{12}-\bar b_1\wedge F_{21})
                         +\alpha\wedge\bar b_1\wedge b_1.
                                                               \tag{1.1}
\]

Here \(F=F^E\). Its global localization, proved in I41, is

\[
 16\pi^2\int d\chi\wedge\mathcal T_\tau
           =2\int\vartheta^2|W^-|^2+\int\vartheta^2 k,
                                                               \tag{1.2}
\] \[
 k=9(|F_A^+|^2-|F_A^-|^2)+s^2/24-|\operatorname{Ric}^{\circ}|^2/2.
\]

The determinant connection on E is 3A. The cutoff equals one near every
zero in the localization convention; thus (1.2) is an identity on the
closed manifold, not an unaccounted boundary formula.

\hypertarget{removing-the-opposite-weyl-variable-without-differentiating-it}{%
\subsection{2. Removing the opposite-Weyl variable without
differentiating
it}\label{removing-the-opposite-weyl-variable-without-differentiating-it}}

Take the compatible ASD basis \(\Omega_1=e^{12}-e^{34},\
\Omega_2=e^{13}+e^{24},\
\Omega_3=e^{14}-e^{23}\), each of squared norm two, and put
\(I_a=c^-(\Omega_a)/2\). Write \(W^-\Omega_b=\sum_a W_{ab}\Omega_a\).
The opposite-Weyl contribution to F is

\[
 F_W=-\tfrac12\sum_{a,b}W_{ab}\Omega_b I_a.
\]

Its entries are

\[
 (F_W)_{22}=\tfrac i2\sum_bW_{1b}\Omega_b,\quad
 (F_W)_{12}=-\tfrac12\sum_bW_{2b}\Omega_b
                         +\tfrac i2\sum_bW_{3b}\Omega_b,
\] \[
 (F_W)_{21}=\tfrac12\sum_bW_{2b}\Omega_b
                         +\tfrac i2\sum_bW_{3b}\Omega_b.
\]

The remainder \(F_0=F-F_W\) obeys \[
 |F_0|\le C(|s|+|\operatorname{Ric}^{\circ}|+|F_A|).
\]

Set \(v_1=\theta,v_2=-\gamma,v_3=-\delta\), and define the matrix and
its symmetric trace-free part by

\[
 K_{ab}=\langle d\vartheta\wedge v_a,\Omega_b\rangle,\qquad
 M=\tfrac12(K+K^{\mathsf T})-\tfrac13(\operatorname{tr}K)I.
\]

Define actual scalar densities j\_chi and q\_chi by

\[
 j_\chi\,d\mu=16\vartheta\,d\vartheta\wedge\theta\wedge\gamma\wedge\delta,
 \qquad
 q_\chi=|d\vartheta|^2(|\theta|^2+|\gamma|^2+|\delta|^2).
\]

These vanish outside T. They are independent of the choice of the
unitary complement to u. Equivalently M is the trace-free symmetric
operator determined by its pairing with the intrinsic linear Weyl term
below; changes of the compatible ASD frame conjugate that operator. No
derivative of a chosen eigenvector of W is used.

Direct substitution in (1.1), using star Omega\_b = -Omega\_b, gives

\[
 4\pi^2(\mathcal T_\tau)_W
       =\tfrac12\sum_{a,b}W_{ab}v_a\wedge\Omega_b,\quad
 4\pi^2(\mathcal T_\tau)_{\mathrm{cubic}}
                                  =-2\theta\wedge\gamma\wedge\delta,
\] \[
 16\pi^2d\chi\wedge(\mathcal T_\tau)_W
                       =4\vartheta\langle M,W^-\rangle\,d\mu.
\]

Consequently define the explicit lower-curvature remainder

\[
 R_\chi=4\int d\chi\wedge\left[
 -\alpha\wedge(F_0)_{22}
 +\tfrac12\{b_1\wedge(F_0)_{12}-\bar b_1\wedge(F_0)_{21}\}\right].
\]

All imaginary coefficients combine into a real integral. Completing the
square in (1.2) now gives the exact global identity

\[
 \boxed{\int\vartheta^2k
 =\int\{j_\chi+2|M|^2-2|\vartheta W^--M|^2\}+R_\chi.}     \tag{2.1}
\]

No second derivative of tau or curvature remains. In particular this
does not assert that the Hessian in the earlier adjoint identity was
pointwise zero; it combines the entire characteristic flux instead. For
every epsilon \textgreater{} 0, Young's inequality gives

\[
 |R_\chi|\le\varepsilon\int q_\chi
 +\frac C\varepsilon\int_T
       (s^2+|\operatorname{Ric}^{\circ}|^2+|F_A|^2).       \tag{2.2}
\]

Indeed its absolute integrand is bounded by C vartheta \textbar d
vartheta\textbar{}
(\textbar theta\textbar+\textbar b\_1\textbar)\textbar F\_0\textbar. All
cutoff derivatives have been retained.

More explicitly, with the Hilbert--Schmidt norm on End(E), orthogonality
of the curvature blocks gives \[
 |F_0|^2=s^2/48+|\operatorname{Ric}^{\circ}|^2/4+(9/2)|F_A|^2.
\] The direct transgression bound is
\(|R_\chi|\le8\int\vartheta|F_0|\sqrt{q_\chi}\). Thus the second term of
(2.2) can be replaced by \[
 \frac1\varepsilon\int_T\vartheta^2
           \{s^2/3+4|\operatorname{Ric}^{\circ}|^2+72|F_A|^2\}.
\] These explicit coefficients are sufficient upper bounds, not a claim
of optimal constants.

\hypertarget{a-topological-curvature-budget-and-the-precise-threshold}{%
\subsection{3. A topological curvature budget and the precise
threshold}\label{a-topological-curvature-budget-and-the-precise-threshold}}

For a fixed sigma \textgreater{} 0 put

\[
 p_{\chi,\sigma}=j_\chi+
                   \frac{4\vartheta^2}{2\vartheta^2+\sigma}|M|^2,
 \quad
 \mathcal S_\sigma=\int(2\vartheta^2+\sigma)
        \left|W^--\frac{2\vartheta M}{2\vartheta^2+\sigma}\right|^2.
\]

At sigma = 0 use \(p_{\chi,0}=j_\chi+2|M|^2\) and
\(\mathcal S_0=2\int|\vartheta W^--M|^2\). The exact refinement of (2.1)
is

\[
 \boxed{\int\vartheta^2k
    =\int p_{\chi,\sigma}+\sigma\int|W^-|^2
                                  -\mathcal S_\sigma+R_\chi.}   \tag{3.1}
\]

For each fixed positive sigma this is the optimal pointwise quadratic
upper bound if W is allowed to range over arbitrary trace-free symmetric
matrices: the square vanishes at W = 2 vartheta M/(2 vartheta squared +
sigma). This algebraic statement is not a claim of geometric
realizability, or of dual attainment under an arbitrary specified global
budget.

Now suppose there were an actual endpoint sequence with beta tending to
infinity, eta = r + D\_beta/V + beta to the power minus two tending to
zero, and the Spin-c class-square condition \(\mathfrak a^2\ge T_+\).
Use normalized integrals throughout this paragraph. From I51 (2.5), rho
tends in L2 to 4/3 and the harmonic-projection remainder J\_H/V tends to
zero. The \emph{actual} Hodge identity is

\[
 4\pi^2p=\int|F_A^+|^2-J_H=\tfrac12\int\rho^2-J_H,
 \qquad p/V\longrightarrow\frac2{9\pi^2}.                \tag{3.2}
\]

This is not obtained by dropping beta times the Weyl energy from the
formula for D\_beta. The bounded ASD fitted model of I51 (2.14a), the
actual torsion-amplitude estimate, and s = -4 yield

\[
 \frac1V\int_T(s^2+|\operatorname{Ric}^{\circ}|^2+|F_A|^2)\to0,
 \qquad
 \frac1V\int\vartheta^2k\to20/3.                         \tag{3.3}
\]

For completeness, the entire low-amplitude set \{m \textless{} b\},
hence also the transition, has relative volume tending to zero; rho is
bounded; the Ricci error tends to zero in normalized L2; and \(F_A^-\)
differs from a bounded field by o(sqrt(V)) in L2. The full averages of
\(|F_A^+|^2,|F_A^-|^2\) tend respectively to 8/9 and 2/9. These prove
both assertions in (3.3), without a pointwise Ricci estimate.

Assume also p \textgreater= A \textgreater{} 0 and chi(M)/A \textless=
B, B \textgreater= 0. Chern--Gauss--Bonnet and (3.2) imply

\[
 \frac{\|W^-\|_2^2}{V}
 \le8\pi^2B\,p/V-\frac23+\frac r2,\qquad
 \limsup\frac{\|W^-\|_2^2}{V}\le
                          \Lambda_B=\frac{16B}{9}-\frac23.     \tag{3.4}
\]

If B \textless{} 3/8 this alone contradicts nonnegativity along such a
sequence. For B \textgreater= 3/8, (2.2)--(3.4) prove the necessary
condition

\[
 \boxed{\liminf_j\left\{
     \fint_j(p_{\chi,\sigma}+\varepsilon q_\chi)
                              +\sigma\Lambda_B\right\}\ge20/3}
                                                               \tag{3.5}
\]

for every \emph{fixed} sigma \textgreater= 0 and epsilon \textgreater{}
0. Therefore any actual subthreshold bound by 20/3 minus a fixed
positive amount, for one such pair, excludes that endpoint sequence.

This asks for a strict subthreshold estimate, not necessarily o(V)
critical energy. The earlier condition that the mixed transition density
is o(V) implies it with sigma = 0, since on the fixed band

\[
 |j_\chi|+|M|^2+q_\chi
           \le C|\langle\tau,\nabla\tau\rangle|^2|\nabla\tau|^2.
\]

No example establishing a strict inclusion between classes of actual
closed solutions is claimed. Most importantly, the original endpoint
budgets do not yet imply this subthreshold estimate.

\hypertarget{the-remaining-cubic-form-can-be-tested-globally}{%
\subsection{4. The remaining cubic form can be tested
globally}\label{the-remaining-cubic-form-can-be-tested-globally}}

The first-derivative density in (2.1) is not an independent topological
constant. Its further Stokes test can also be made exact. The three-form
theta wedge gamma wedge delta is invariant under a change of complement
to u; after multiplication by vartheta squared it extends smoothly by
zero through the zero set of tau. Thus

\[
 \int j_\chi=-8\int\vartheta^2d(\theta\wedge\gamma\wedge\delta).
                                                               \tag{4.1}
\]

Write \(\alpha_2=i\theta_2\),
\(f_{11}=-iF_{11},\ p_{21}=\operatorname{Re}F_{21},\
q_{21}=\operatorname{Im}F_{21}\). The full U(2) structure equations give

\[
 d\theta=f_{11}+2\gamma\wedge\delta,\quad
 d\gamma=p_{21}+(\theta_2-\theta)\wedge\delta,\quad
 d\delta=q_{21}-(\theta_2-\theta)\wedge\gamma.
\]

The complement connection terms cancel, leaving

\[
 d(\theta\wedge\gamma\wedge\delta)
 =f_{11}\wedge\gamma\wedge\delta-\theta\wedge p_{21}\wedge\delta
                                   +\theta\wedge\gamma\wedge q_{21}.
                                                               \tag{4.2}
\]

Let L be the matrix whose three rows pair
\((\gamma\wedge\delta,\theta\wedge\delta,-\theta\wedge\gamma)\) with
Omega\_b, and let L\_0 be its symmetric trace-free part. The W-minus
part of (4.2) is \(\tfrac12\langle W^-,L_0\rangle d\mu\). Define R\_bulk
to be -8 times the integral of vartheta squared times the right side of
(4.2) with F replaced by F\_0. Then

\[
 \int j_\chi=-4\int\vartheta^2\langle W^-,L_0\rangle+R_{\mathrm{bulk}}.
\]

Put H = M-vartheta L\_0. A second exact completion gives

\[
 \boxed{\int\vartheta^2k
     =2\int|H|^2-2\int|\vartheta W^--H|^2
                                     +R_\chi+R_{\mathrm{bulk}}.} \tag{4.3}
\]

Here R\_bulk is supported on the entire nonzero-amplitude region, not
merely T. Its calibrated scalar and determinant terms need not be small.
Equation (4.3) therefore does not erase the critical payment: it moves
it to curvature paired with squared derivatives of the unit torsion in
the bulk. Declaring R\_bulk to be o(V) on the basis of (3.3), which only
concerns T, would be incorrect.

The two exact forms (2.1) and (4.3) are alternative global tests of the
same actual field. A proof of (3.5)'s strict reverse bound, or a
controlled signed cancellation in (4.3), is still required to settle
O10. Neither identity supplies a global escape sequence.

\hypertarget{retaining-the-non-small-calibrated-bulk-curvature}{%
\subsection{5. Retaining the non-small calibrated bulk
curvature}\label{retaining-the-non-small-calibrated-bulk-curvature}}

The bulk remainder in (4.3) can be expanded instead of being estimated
as an unspecified error. In the same actual compatible frame put \[
 F_*=i(\Omega_++\tfrac12\Omega_1)I
                               +\tfrac16\sum_a\Omega_a I_a.
\] Here Omega\_+ is the unit-length-squared-two direction of the actual
positive spinor. The actual amplitude and ASD fit imply \[
 \fint_{\{m\ge a\}}|F_0-F_*|^2\longrightarrow0.
\] The integral is normalized by the full volume. This follows from rho
tending to 4/3, the bounded fitted ASD curvature coefficient tending to
1/3 in the tau-adapted direction, and the vanishing Ricci error. It does
not provide smallness after multiplication by an uncontrolled quadratic
derivative.

In (4.2), the entries of F\_* are \[
 f_{11}^*=\Omega_++\tfrac23\Omega_1,\quad
 p_{21}^*=-\Omega_2/6,\quad q_{21}^*=-\Omega_3/6.
\] Their contribution to R\_bulk is the integral of vartheta squared
times \[
 q_*=-8\langle\gamma\wedge\delta,\Omega_+\rangle
       +\tfrac{16}{3}\langle\gamma\wedge\delta,\Omega_1\rangle
       +\tfrac43\langle\theta\wedge\delta,\Omega_2\rangle
       -\tfrac43\langle\theta\wedge\gamma,\Omega_3\rangle.
\]

For \(u=\tau/\sqrt m\), the exact quadratic identity is \[
 \boxed{q_*=\tfrac{20}{3}|\nabla u|^2-2|\theta|^2
       -8|(\nabla u)^{1,0}|^2-\tfrac23|D^-u|^2
                                    -2|\delta-J_-\gamma|^2.}   \tag{5.1}
\] Here J\_- sends e\_1 to e\_2 and e\_3 to -e\_4, and the (1,0)
projection uses the actual positive-spinor almost-complex structure. The
tensor norm includes the factor one half for evaluation on Z\_1,Z\_2.
One verifies (5.1) from \[
 2|(\nabla u)^{1,0}|^2-|\nabla u|^2
                               =2\langle\gamma\wedge\delta,\Omega_+\rangle,
\] \[
 |D^-u|^2=|\nabla u|^2-2\operatorname{tr}L,\qquad
 2\langle\gamma\wedge\delta,\Omega_1\rangle
             =|\gamma|^2+|\delta|^2-|\delta-J_-\gamma|^2.
\]

The residual calibration error is exactly (4.2) with F\_0-F\_*, and its
absolute value is at most \[
 C\int\vartheta^2|F_0-F_*|
       (|\gamma\wedge\delta|+|\theta\wedge\gamma|+|\theta\wedge\delta|).
                                                               \tag{5.2}
\] The curvature factor is small in normalized L2; the other factor is
quadratic in the actual unit derivative and has no inherited L2 bound.
Also \[
 D^-u=m^{-1/2}D^-\tau-\frac1{2m}c^-(dm)u.
\] Hence the radial term cannot be discarded using weak Dirac control of
tau alone. Equations (5.1)--(5.2) retain both the non-small calibrated
potential and the exact critical error. No strict bound below 20V/3 has
been established from them.

\hypertarget{exact-cancellation-of-the-central-and-scalar-curvature-payments}{%
\subsection{6. Exact cancellation of the central and scalar curvature
payments}\label{exact-cancellation-of-the-central-and-scalar-curvature-payments}}

The determinant part of (5.2) need not be bounded in absolute value.
There is a further global cancellation. It removes that derivative
payment exactly, but also removes the positive characteristic budget;
both changes must be retained.

Write h = vartheta squared = 1 - chi and decompose the actual curvature
on E as \[
 F^E=F_W+F_s+F_R+i c I,\qquad
 c=\tfrac32(-iF_A),\quad dc=0,\quad
 F_s=\tfrac16\sum_a\Omega_a I_a.
\] The scalar is exactly -4 in the normalization used throughout this
calculation. The Ricci block has the form \[
 F_R=\begin{pmatrix}ir&-p+iq\\p+iq&-ir\end{pmatrix},
 \qquad r,p,q\in\Lambda^+,
 \qquad |F_R|_{\mathrm{HS}}^2=\tfrac14|\operatorname{Ric}^{\circ}|^2.
\] These letters r,p,q are real two-forms in this section, not a radial
coordinate, a monopole budget or an error norm. Set \[
 x_b=\tfrac12\langle c^-,\Omega_b\rangle,\qquad
 C=\operatorname{STF}(e_1\otimes x),\qquad
 f_{11}=-iF^E_{11}
   =c-\tfrac12\sum_bW_{1b}^-\Omega_b+\Omega_1/6+r.
\] The first structure equation is
\(d\theta=f_{11}+2\gamma\wedge\delta\).

For any curvature block G define the real three- and four-forms \[
 S(G)=-\alpha\wedge G_{22}
       +\tfrac12(b_1\wedge G_{12}-\overline b_1\wedge G_{21}),
\] \[
 B(G)=(-iG_{11})\wedge\gamma\wedge\delta
       -\theta\wedge\operatorname{Re}G_{21}\wedge\delta
       +\theta\wedge\gamma\wedge\operatorname{Im}G_{21}.
\] Only the displayed actual blocks, for which these expressions are
real, are used. Put \(R_\chi(G)=4\int d\chi\wedge S(G)\) and
\(R_b(G)=-8\int hB(G)\).

\hypertarget{the-central-block}{%
\subsubsection{6.1 The central block}\label{the-central-block}}

Here S(icI) = theta wedge c and B(icI) = c wedge gamma wedge delta. The
one-form theta is intrinsically determined by the actual unit section u;
multiplication by h extends it smoothly by zero near tau = 0. Since dc =
0, Stokes gives \[
 \int d(h\theta\wedge c)=0,
 \qquad
 \boxed{R_\chi(icI)+R_b(icI)=4\int h c\wedge f_{11}.}     \tag{6.1}
\] Indeed the two gamma wedge delta terms cancel exactly. No norm of the
derivative of c or any smallness of c has been used.

\hypertarget{the-scalar-block-and-the-full-cutoff-term}{%
\subsubsection{6.2 The scalar block and the full cutoff
term}\label{the-scalar-block-and-the-full-cutoff-term}}

All spinor pairings below are the real parts of Hermitian pairings,
conjugate-linear in the first entry. Define \[
 X_h=\operatorname{Re}\int
        \langle D^-u,c^-(dh)u\rangle.
\] The exact weighted Lichnerowicz identity is \[
 \int h|D^-u|^2+X_h
      =\int h|\nabla u|^2-\int h(1+2x_1).                \tag{6.2}
\] The scalar potential is -1, and the central ASD curvature acts on u =
e\_1 by -2x\_1. The Ricci block is self-dual and gives no Dirac
potential on this negative spinor. There is no gradient-cutoff term from
the connection Laplacian because \(|u|=1\), so its pairing with the
derivative of u has zero real part.

With the Clifford matrices fixed in the preceding actual-prolongation
proof and \(\nabla u=(i\theta,\gamma+i\delta)\), finite Clifford algebra
gives \[
 X_h=\int\sum_a\langle dh\wedge v_a,\Omega_a\rangle,
                     \qquad v=(\theta,-\gamma,-\delta).
\] In particular the cutoff term is not omitted. Direct evaluation
yields \[
 S(F_s)=-\tfrac16\sum_a v_a\wedge\Omega_a,
 \quad R_\chi(F_s)=-\tfrac23 X_h,
 \quad R_b(F_s)=\tfrac23\int h(|\nabla u|^2-|D^-u|^2).
\] Using (6.2), \[
 \boxed{R_\chi(F_s)+R_b(F_s)=\tfrac23\int h(1+2x_1).}    \tag{6.3}
\] This step does not assume D\^{}-tau = 0 or discard the radial part of
D\^{}-u.

Equivalently, the entire scalar contribution is
\(\tfrac23(\|\nabla(\vartheta u)\|_2^2 -\|D^-(\vartheta u)\|_2^2)\).
Apply Lichnerowicz to the globally smooth section vartheta u, which
vanishes in a neighborhood of the zero set. The two cutoff-gradient
squares cancel. This gives the same formula without requiring u itself
to extend through that set.

\hypertarget{the-remaining-ricci-interaction-and-the-lost-positive-budget}{%
\subsubsection{6.3 The remaining Ricci interaction and the lost positive
budget}\label{the-remaining-ricci-interaction-and-the-lost-positive-budget}}

Since the Omega\_a are anti-self-dual, \[
 4c\wedge f_{11}
 =4c\wedge c+
     (4\langle W^-,C\rangle-\tfrac43x_1)d\mu
       +4c\wedge r.
\] The x\_1 term cancels with (6.3). Moreover the characteristic density
in Section 1 is exactly \[
                 k=4\frac{c\wedge c}{d\mu}
                      +\tfrac23-\tfrac12|\operatorname{Ric}^{\circ}|^2.
\] Substitute these two equalities into (4.3), and put
\(\widetilde H=M-\vartheta L_0+\vartheta C\). The result is \[
\boxed{\begin{aligned}
 0={}&2\int|\widetilde H|^2
       -2\int|\vartheta W^--\widetilde H|^2
       +\mathcal R_R\\
     &+4\int h c\wedge r
       +\tfrac12\int h|\operatorname{Ric}^{\circ}|^2,
\end{aligned}}                                                    \tag{6.4}
\] where the sole derivative-curvature payment is \[
\begin{aligned}
 \mathcal R_R={}&4\int dh\wedge
          (\theta\wedge r+\gamma\wedge p+\delta\wedge q)\\
 &-8\int h\{r\wedge\gamma\wedge\delta
       -\theta\wedge p\wedge\delta+\theta\wedge\gamma\wedge q\}.
\end{aligned}                                                     \tag{6.5}
\] All wedge products are actual global contractions of the adapted
connection expressions; frame changes preserve the paired integral. The
last two terms in (6.4) are o(V) along the endpoint sequence: c has
bounded normalized L2 norm and Ricci trace-free error tends to zero in
that norm. This observation does not apply to (6.5), which contains
squared derivatives of the unit field.

Thus the determinant-error product in (5.2) has been removed globally,
not estimated. The whole positive 20V/3 budget has also disappeared into
the central and scalar identities. Even a proof that
\(\mathcal R_R=o(V)\) would give equality of the two square norms in
(6.4), not a contradiction. The missing statement must be a strict
signed separation compatible with the actual cubic provenance and
bracket, or a joint estimate of that separation and (6.5). These
identities hold for a general unit section on the cutoff support; their
validity alone does not exploit all of those additional actual field
equations. No unrestricted endpoint gap follows.

\hypertarget{a-single-intrinsic-mixed-curvature-pairing}{%
\subsubsection{6.4 A single intrinsic mixed-curvature
pairing}\label{a-single-intrinsic-mixed-curvature-pairing}}

The remaining terms can be combined without taking absolute values. Let
Psi\_b be the compatible positive two-form basis with squared norm two,
and write \[
 F_R=-\tfrac12\sum_{a,b}\mathsf R_{ab}\Psi_b I_a,
 \qquad |\mathsf R|^2=\tfrac14|\operatorname{Ric}^{\circ}|^2.
\] Use the full three-by-three matrices \[
 K^+_{ab}=\langle d\vartheta\wedge v_a,\Psi_b\rangle,
 \quad L^+_{ab}=\langle A_a,\Psi_b\rangle,
 \quad A=(\gamma\wedge\delta,\theta\wedge\delta,-\theta\wedge\gamma),
\] \[
 y_b=\tfrac12\langle c^+,\Psi_b\rangle,
 \qquad H^+=K^+-\vartheta L^++\vartheta(e_1\otimes y).
\] For the actual monopole \(c^+=(3\rho/4)\Omega_+\), so
\(y=(3\rho/4,0,0)\) in this frame. In contrast with \(\widetilde H\),
H\^{}+ is \textbf{not} projected to its symmetric trace-free part. The
mixed curvature is a map between the two different handed two-form
bundles.

Substitution into (6.5), using self-duality for the positive wedge
products, gives exactly \[
 \mathcal R_R+4\int h c\wedge r
                         =-4\int\vartheta\langle\mathsf R,H^+\rangle.
\] Hence the final identity is \[
 \boxed{\|\vartheta W^--\widetilde H\|_2^2
                 -\|\widetilde H\|_2^2
       =\|\vartheta\mathsf R\|_2^2
                    -2\int\vartheta\langle\mathsf R,H^+\rangle.} \tag{6.6}
\] This retains the sign, the cutoff and every curvature block. H\^{}+
contains quadratic unit-field derivatives; no uniform normalized L2
bound for it follows from the stated budgets. Moreover vanishing of the
right side would still force only a zero square difference. The next
sufficient estimate must establish a strictly signed separation using
the cubic field equations, jointly with the mixed pairing. Equation
(6.6) is the current exact continuation point, not an added field
equation or a completed endpoint theorem.

Testing the unnormalized actual tau does not restore a discarded radial
coercivity in this channel. Put \(\lambda=dm/(2m)\) and use the smooth
weight h/m, extended by zero. Substituting
\(\nabla\tau=\sqrt m(\nabla u+\lambda\otimes u)\),
\(D^-\tau=\sqrt m(D^-u+c^-(\lambda)u)\) into weighted Lichnerowicz adds
the identical term \(-\int h|\lambda|^2+\int\langle dh,\lambda\rangle\)
to both sides of (6.2). It cancels exactly. Thus a repeated scalar test
of tau is not a missing source of a positive length-gradient term; the
additional cubic and bracket information must enter the signed curvature
identity itself.

\hypertarget{the-cubic-actual-core-does-not-supply-the-missing-signed-separation}{%
\subsection{7. The cubic actual core does not supply the missing signed
separation}\label{the-cubic-actual-core-does-not-supply-the-missing-signed-separation}}

The next test uses the actual fixed-hyperbolic, weight-two monopole
solutions of \href{20260927_conformal_separation_and_zero_end.md}{I41,
Section 15.1}, not independently assigned spinor jets. Normalize their
simple torsion zero so that, on shrinking mass shells, the physical
field has leading term

\[
 \tau=(\overline z_1,\overline z_2)/\ell,
 \qquad m=|\tau|^2=r^2/\ell^2,
 \qquad \rho\longrightarrow\rho_*=4/3.
 \tag{7.1}
\]

Fixed constant factors in the original construction scale are absorbed
into ell. All following limits concern a fixed interval \(0<a<b<128/81\)
in m and the C1 rescaled expansions of that actual construction. The
metric is exactly hyperbolic: \(\operatorname{Ric}=-g\), \(W^-=0\) and
\(\mathsf R=0\). Thus the two sides of (6.6) vanish identically, even
though the actual cubic bracket is present. Its validity alone cannot
supply a strictly positive local square separation.

There is a quantitative version which retains the full cutoff. Set
\(\vartheta=\vartheta_0(m)\), with values zero and one near the two ends
of the shell. At \(x=(r,0,0,0)\), the leading unit-field forms are

\[
 \theta=-e^2/r,\qquad \gamma=e^3/r,\qquad\delta=-e^4/r.
\]

Using the compatible two-form frames of Section 6, direct substitution
gives the frame-invariant leading norms

\[
 |\widetilde H|^2=\frac{8}{3r^4}(r\vartheta_r+\vartheta)^2,
 \qquad |H^+|^2=\frac{3}{r^4}(r\vartheta_r-\vartheta)^2.
\]

The central bounded terms contribute o(1) to their shell integrals.
Since \(r\vartheta_r=2m\vartheta_0'\), the limits are

\[
 E_-:=\lim\int_{a<m<b}|\widetilde H|^2
   =\frac{8\pi^2}{3}\int_a^b(2m\vartheta_0'+\vartheta_0)^2\frac{dm}{m},
\] \[
 E_+:=\lim\int_{a<m<b}|H^+|^2
   =3\pi^2\int_a^b(2m\vartheta_0'-\vartheta_0)^2\frac{dm}{m}.
 \tag{7.2}
\]

Subtracting and using the endpoint values proves the exact relation and
sharp cutoff infima

\[
 E_- -\frac89 E_+=\frac{32\pi^2}{3},\qquad
 \inf E_-=\frac{32\pi^2b}{3(b-a)},\qquad
 \inf E_+=\frac{12\pi^2a}{b-a}. \tag{7.3}
\]

Indeed, writing \(\vartheta_0=\sqrt m\,v\) makes the E-plus integral
\(12\pi^2\int_a^b m^2(v')^2dm\), with \(v(a)=0\), \(v(b)=1/\sqrt b\).
Weighted Cauchy gives its infimum; the minimizing profile is
\(\vartheta_0(m)=(m-a)\sqrt b/[(b-a)\sqrt m]\). Smooth flattening near
the endpoints approaches it. This is a shell cost calculation, not a
lower bound for the original Einstein defect.

To test whether these shells are already excluded by the actual
calibrated field fit, use I51's bounded model. Here \(S=m_*+m+o(1)\),
\(m_*=128/81\), the actual ASD determinant curvature f is O(ell), and

\[
 K=\frac{2m}{3(m_*+m)}\Omega_-+o(1),\qquad |\Omega_-|^2=2.
\]

Consequently

\[
 \int_{a<m<b}|f-K|^2
  =\frac{8\pi^2\ell^4}{9}\int_a^b\frac{m^3}{(m_*+m)^2}\,dm
     +o(\ell^4),
\] \[
 \int_{a<m<b}|\nabla\tau|^2
       =2\pi^2\ell^2(b^2-a^2)+o(\ell^2). \tag{7.4}
\]

Thus the actual fit pays volume scale and the quadratic derivative
budget pays scale squared, while the two critical norms in (7.2) retain
nonzero limits. These local actual solutions remain compatible with
those subcritical budgets under shrinking. Their complete boundary
matching and global conformal/weight supply are not proved. They neither
close O10 nor constitute a closed escape sequence.

The useful consequence for the next proof is precise: a strict
separation must use information carried by the calibrated exterior and
its matching with these cores, or another full-system identity. The
actual cubic bracket on a core, by itself, does not restore the positive
budget cancelled in Section 6. No universal estimate for the mixed Ricci
term follows from the cutoff infima (7.3).

\end{document}
